% Общие параметры блок-схем (ГОСТ 19.701): единая ширина блоков
\newcommand{\flowchartBlockWidth}{6.8cm}
\newcommand{\flowchartBlockHeight}{0.85cm}
\newcommand{\flowchartRightShift}{6.3cm}      % отступ правой колонки (рис. 2.2--2.8)
\newcommand{\flowchartRightShiftSmall}{5.0cm} % компактная схема calculate_ts
\newcommand{\flowchartUpdateVGap}{1.28cm}     % верт. зазор для рис. 2.10
\newcommand{\flowchartUpdateMidExtra}{-0.4cm}   % pro3→pro4, pro4→pro5 (рис. 2.10)
\newcommand{\flowchartTinyVExtra}{-0.15cm}    % микро-добавка (рис. 2.3)
\newcommand{\flowchartDecExpExtra}{-0.45cm}   % dec1 → expected (рис. 2.2)

\tikzset{
  fc base/.style={font=\footnotesize, inner sep=4pt},
  fc arrow/.style={thick,->,>=stealth},
  fc startstop/.style={draw, rectangle, rounded corners=10pt,
    text width=\flowchartBlockWidth, minimum height=\flowchartBlockHeight,
    align=center, fill=white},
  fc process/.style={draw, rectangle,
    text width=\flowchartBlockWidth, minimum height=\flowchartBlockHeight,
    align=center, fill=white},
  fc predefined/.style={draw, rectangle,
    text width=\flowchartBlockWidth, minimum height=\flowchartBlockHeight,
    align=center, fill=white,
    path picture={
      \draw ([xshift=4pt]path picture bounding box.north west) --
            ([xshift=4pt]path picture bounding box.south west);
      \draw ([xshift=-4pt]path picture bounding box.north east) --
            ([xshift=-4pt]path picture bounding box.south east);
    }},
  fc io/.style={draw, trapezium, trapezium left angle=75, trapezium right angle=105,
    trapezium stretches=true,
    text width=\flowchartBlockWidth, minimum height=\flowchartBlockHeight,
    align=center, fill=white},
  fc decision/.style={draw, diamond, aspect=3.2,
    text width=5.0cm, inner sep=1pt,
    minimum height=\flowchartBlockHeight, align=center, fill=white},
}

\newcommand{\flowchartCalculateTs}{%
\begin{figure}[H]
\centering
\begin{tikzpicture}[
    fc base,
    node distance=1.5cm,
    startstop/.style={fc startstop},
    process/.style={fc process},
    io/.style={fc io},
    decision/.style={fc decision},
    arrow/.style={fc arrow}
]
\node (start)  [startstop] {Начало};
\node (pro1)   [process,  below of=start] {$PTS, FGA, FTA := \mathrm{SUM}(\ldots)$};
\node (pro2)   [process,  below of=pro1]
  {$D = 2 \cdot (\mathrm{FGA} + 0{,}44 \cdot \mathrm{FTA})$};
\node (dec1)   [decision, below of=pro2, yshift=-0.5cm] {$D = 0$?};
\node (pro3a)  [io,  right of=dec1, xshift=\flowchartRightShiftSmall] {Вывод:\\ значение не определено};
\node (pro3b)  [io,  below of=dec1, yshift=-0.5cm]
  {Возврат $\mathrm{PTS} / D$};
\node (stop)   [startstop, below of=pro3b] {Конец};

\draw [arrow] (start)  -- (pro1);
\draw [arrow] (pro1)   -- (pro2);
\draw [arrow] (pro2)   -- (dec1);
\draw [arrow] (dec1)   -- node[anchor=south] {Да}  (pro3a);
\draw [arrow] (dec1)   -- node[anchor=east]  {Нет} (pro3b);
\draw [arrow] (pro3b)  -- (stop);
\draw [arrow] (pro3a)  |- (stop);
\end{tikzpicture}
\caption{Схема алгоритма вычисления функции \texttt{calculate\_ts}}
\label{fig:algo-ts}
\end{figure}
}

\newcommand{\flowchartCalculateEfg}{%
\begin{figure}[H]
\centering
\begin{tikzpicture}[
    fc base,
    node distance=1.8cm,
    startstop/.style={fc startstop},
    process/.style={fc process},
    io/.style={fc io},
    decision/.style={fc decision},
    arrow/.style={fc arrow}
]
\node (start)  [startstop] {Начало};
\node (agg)    [process,  below of=start]
  {$FGM, FG3M, FGA := \mathrm{SUM}(\ldots)$ \\ из \texttt{game\_player\_stats}};
\node (dec1)   [decision, below of=agg] {$\mathrm{FGA} = 0$?};
\node (null1)  [io,  right of=dec1, xshift=\flowchartRightShift] {Вывод:\\ значение не определено};
\node (ret)    [io,  below of=dec1]
  {Возврат $\displaystyle\frac{\mathrm{FGM}+0{,}5\cdot\mathrm{FG3M}}{\mathrm{FGA}}$};
\node (stop)   [startstop, below of=ret] {Конец};

\draw [arrow] (start) -- (agg);
\draw [arrow] (agg)   -- (dec1);
\draw [arrow] (dec1)  -- node[above] {Да} (null1);
\draw [arrow] (dec1)  -- node[right] {Нет} (ret);
\draw [arrow] (ret)   -- (stop);
\draw [arrow] (null1.east) -- ++(0.6,0) |- (stop);
\end{tikzpicture}
\caption{Схема алгоритма вычисления функции \texttt{calculate\_efg}}
\label{fig:algo-efg}
\end{figure}
}

\newcommand{\flowchartCalculateUsg}{%
\begin{figure}[H]
\centering
\begin{tikzpicture}[
    fc base,
    node distance=1.8cm,
    startstop/.style={fc startstop},
    process/.style={fc process},
    io/.style={fc io},
    decision/.style={fc decision},
    arrow/.style={fc arrow}
]
\node (start)  [startstop] {Начало};
\node (tid)    [process,  below of=start]
  {Определить team\_id \\ (параметр или последняя игра)};
\node (dec0)   [decision, below of=tid, yshift=-0.2cm] {team\_id IS NULL?};
\node (null0)  [io,  right of=dec0, xshift=\flowchartRightShift] {Вывод:\\ значение не определено};
\node (aggP)   [process,  below of=dec0, yshift=-0.2cm]
  {$pFGA, pFTA, pTOV, pMP :=$ \\ $\mathrm{SUM}(\ldots)$ игрока};
\node (dec1)   [decision, below of=aggP, yshift=-0.2cm] {$pMP < 1$?};
\node (null1)  [io,  right of=dec1, xshift=\flowchartRightShift] {Вывод:\\ значение не определено};
\node (aggT)   [process,  below of=dec1, yshift=-0.2cm]
  {$tmFGA, tmFTA, tmTOV, tmMP :=$ \\ $\mathrm{SUM}(\ldots)$ команды};
\node (denom)  [process,  below of=aggT, yshift=-0.2cm]
  {$denom := pMP \cdot (tmFGA + 0{,}44 \cdot tmFTA + tmTOV)$};
\node (dec2)   [decision, below of=denom, yshift=-0.2cm] {denom $= 0$?};
\node (null2)  [io,  right of=dec2, xshift=\flowchartRightShift] {Вывод:\\ значение не определено};
\node (ret)    [io,  below of=dec2, yshift=-0.2cm]
  {Возврат $\displaystyle\frac{(pFGA+0{,}44\cdot pFTA+pTOV)\cdot\frac{tmMP}{5}}{denom}$};
\node (stop)   [startstop, below of=ret] {Конец};

\draw [arrow] (start) -- (tid);
\draw [arrow] (tid)   -- (dec0);
\draw [arrow] (dec0)  -- node[above] {Да} (null0);
\draw [arrow] (dec0)  -- node[right] {Нет} (aggP);
\draw [arrow] (aggP)  -- (dec1);
\draw [arrow] (dec1)  -- node[above] {Да} (null1);
\draw [arrow] (dec1)  -- node[right] {Нет} (aggT);
\draw [arrow] (aggT)  -- (denom);
\draw [arrow] (denom) -- (dec2);
\draw [arrow] (dec2)  -- node[above] {Да} (null2);
\draw [arrow] (dec2)  -- node[right] {Нет} (ret);
\draw [arrow] (ret)   -- (stop);

\draw [arrow] (null0.east) -- ++(0.5,0) coordinate (bus) |- (stop);
\draw [arrow] (null1.east) -- (null1.east -| bus);
\draw [arrow] (null2.east) -- (null2.east -| bus);
\end{tikzpicture}
\caption{Схема алгоритма вычисления функции \texttt{calculate\_usg}}
\label{fig:algo-usg}
\end{figure}
}

\newcommand{\flowchartCalculatePer}{%
\begin{figure}[H]
\centering
\begin{tikzpicture}[
    fc base,
    node distance=1.8cm,
    startstop/.style={fc startstop},
    process/.style={fc process},
    predefined/.style={fc predefined},
    io/.style={fc io},
    decision/.style={fc decision},
    arrow/.style={fc arrow}
]
\node (start) [startstop] {Начало};
\node (agg)   [process,  below of=start]
  {$PTS, TRB, AST, STL, BLK,$ \\ $TOV, PF, FGM/FGA, FTM/FTA, MP := \mathrm{SUM}(\ldots)$};
\node (dec1)  [decision, below of=agg] {$\mathrm{MP} < 50$?};
\node (null1) [io,  right of=dec1, xshift=\flowchartRightShift] {Вывод:\\ значение не определено};
\node (uper)  [process,  below of=dec1]
  {$\mathrm{uPER} = (\mathrm{PTS}+\mathrm{TRB}+\ldots - \mathrm{TOV}-\mathrm{PF}-\Delta\mathrm{FG}-0{,}44\cdot\Delta\mathrm{FT}) \cdot \tfrac{48}{\mathrm{MP}}$};
\node (lg)    [predefined,  below of=uper]
  {Вычисление lg\_aPER: \\ AVG(uPER) по игрокам с MP $> 100$};
\node (dec2)  [decision, below of=lg] {lg\_aPER $= 0$ или NULL?};
\node (fix)   [process,  right of=dec2, xshift=\flowchartRightShift] {lg\_aPER $= 15{,}0$};
\node (ret)   [io,  below of=dec2]
  {Возврат $\mathrm{uPER} \cdot 15\,/\,\mathrm{lg\_aPER}$};
\node (stop)  [startstop, below of=ret] {Конец};

\draw [arrow] (start) -- (agg);
\draw [arrow] (agg)   -- (dec1);
\draw [arrow] (dec1)  -- node[above] {Да} (null1);
\draw [arrow] (dec1)  -- node[right] {Нет} (uper);
\draw [arrow] (uper)  -- (lg);
\draw [arrow] (lg)    -- (dec2);
\draw [arrow] (dec2)  -- node[above] {Да} (fix);
\draw [arrow] (dec2)  -- node[right] {Нет} (ret);
\draw [arrow] (fix)   |- (ret);
\draw [arrow] (ret)   -- (stop);
\draw [arrow] (null1.east) -- ++(1.0,0) |- (stop);
\end{tikzpicture}
\caption{Схема алгоритма вычисления функции \texttt{calculate\_per}}
\label{fig:algo-per}
\end{figure}
}

\newcommand{\flowchartCalculateBpm}{%
\begin{figure}[H]
\centering
\begin{tikzpicture}[
    fc base,
    node distance=1.8cm,
    startstop/.style={fc startstop},
    process/.style={fc process},
    predefined/.style={fc predefined},
    io/.style={fc io},
    decision/.style={fc decision},
    arrow/.style={fc arrow}
]
\node (start) [startstop] {Начало};
\node (agg)   [process,  below of=start]
  {$PTS, TRB, AST, STL,$ \\ $BLK, TOV, PF, MP := \mathrm{SUM}(\ldots)$};
\node (dec1)  [decision, below of=agg] {$\mathrm{MP} < 50$?};
\node (null1) [io,  right of=dec1, xshift=\flowchartRightShift] {Вывод:\\ значение не определено};
\node (norm)  [process,  below of=dec1]
  {Нормировка на 100 мин: $X_{100} = X \cdot 100\,/\,\mathrm{MP}$};
\node (raw)   [process,  below of=norm]
  {raw $= 0{,}123\cdot\mathrm{PTS}_{100}+0{,}234\cdot\mathrm{TRB}_{100}+0{,}689\cdot\mathrm{AST}_{100}$ \\ $+\,0{,}445\cdot\mathrm{STL}_{100}+0{,}407\cdot\mathrm{BLK}_{100}-0{,}605\cdot\mathrm{TOV}_{100}-0{,}086\cdot\mathrm{PF}_{100}$};
\node (lg)    [predefined,  below of=raw]
  {Вычисление lg\_avg: \\ AVG(raw) по игрокам с MP $> 100$};
\node (ret)   [io,  below of=lg] {Возврат raw $-$ lg\_avg};
\node (stop)  [startstop, below of=ret] {Конец};

\draw [arrow] (start) -- (agg);
\draw [arrow] (agg)   -- (dec1);
\draw [arrow] (dec1)  -- node[above] {Да} (null1);
\draw [arrow] (dec1)  -- node[right] {Нет} (norm);
\draw [arrow] (norm)  -- (raw);
\draw [arrow] (raw)   -- (lg);
\draw [arrow] (lg)    -- (ret);
\draw [arrow] (ret)   -- (stop);
\draw [arrow] (null1.east) -- ++(0.6,0) |- (stop);
\end{tikzpicture}
\caption{Схема алгоритма вычисления функции \texttt{calculate\_bpm}}
\label{fig:algo-bpm}
\end{figure}
}

\newcommand{\flowchartTriggerScore}{%
\begin{figure}[H]
\centering
\begin{tikzpicture}[
    fc base,
    node distance=1.8cm,
    startstop/.style={fc startstop},
    process/.style={fc process},
    predefined/.style={fc predefined},
    io/.style={fc io},
    decision/.style={fc decision},
    arrow/.style={fc arrow}
]
\node (start) [startstop] {Начало};
\node (get)   [io,  below of=start]
  {Получить status, home/away\_id, \\ счёт из \texttt{games} по game\_id};
\node (dec1)  [decision, below of=get] {status $\neq$ 'Finished' \\ или счёт NULL?};
\node (skip)  [io,  right of=dec1, xshift=\flowchartRightShift] {Возврат NEW/OLD};
\node (exp)   [process,  below of=dec1, yshift=\flowchartDecExpExtra]
  {$expected :=$ счёт команды \\ (\texttt{home\_score} или \texttt{away\_score})};
\node (loop)  [predefined,  below of=exp]
  {Для каждой команды в матче: \\ $teamPoints := \mathrm{SUM(points)}$};
\node (dec2)  [decision, below of=loop] {$teamPoints \neq expected$?};
\node (exc)   [io,  right of=dec2, xshift=\flowchartRightShift] {RAISE EXCEPTION};
\node (ret)   [io,  below of=dec2] {Возврат NEW/OLD};
\node (stop)  [startstop, below of=ret] {Конец};

\draw [arrow] (start) -- (get);
\draw [arrow] (get)   -- (dec1);
\draw [arrow] (dec1)  -- node[above] {Да} (skip);
\draw [arrow] (dec1)  -- node[right] {Нет} (exp);
\draw [arrow] (exp)   -- (loop);
\draw [arrow] (loop)  -- (dec2);
\draw [arrow] (dec2)  -- node[above] {Да} (exc);
\draw [arrow] (dec2)  -- node[right] {Нет} (ret);
\draw [arrow] (ret)   -- (stop);

\draw [arrow] (skip.east) -- ++(0.6,0) coordinate (bus) |- (stop);
\draw [arrow] (exc.east) -- (exc.east -| bus);
\end{tikzpicture}
\caption{Схема алгоритма функции-триггера \texttt{trg\_validate\_team\_score\_row}}
\label{fig:algo-trg-score}
\end{figure}
}

\newcommand{\flowchartTriggerSeasonStats}{%
\begin{figure}[H]
\centering
\begin{tikzpicture}[
    fc base,
    node distance=1.8cm,
    startstop/.style={fc startstop},
    process/.style={fc process},
    predefined/.style={fc predefined},
    io/.style={fc io},
    decision/.style={fc decision},
    arrow/.style={fc arrow}
]
\node (start) [startstop] {Начало};
\node (get)   [io,  below of=start]
  {Получить список уникальных season\_id \\ из новых строк};
\node (dec1)  [decision, below of=get] {Список season\_id пуст?};
\node (null1) [io,  right of=dec1, xshift=\flowchartRightShift] {Вывод:\\ значение не определено};
\node (call)  [predefined,  below of=dec1, yshift=\flowchartTinyVExtra]
  {Для каждого season\_id: \\ \texttt{CALL update\_season\_stats(season\_id)}};
\node (ret)   [io,  below of=call, yshift=\flowchartTinyVExtra] {Вывод:\\ значение не определено};
\node (stop)  [startstop, below of=ret] {Конец};

\draw [arrow] (start) -- (get);
\draw [arrow] (get)   -- (dec1);
\draw [arrow] (dec1)  -- node[above] {Да} (null1);
\draw [arrow] (dec1)  -- node[right] {Нет} (call);
\draw [arrow] (call)  -- (ret);
\draw [arrow] (ret)   -- (stop);
\draw [arrow] (null1.east) -- ++(0.6,0) |- (stop);
\end{tikzpicture}
\caption{Схема алгоритма функции-триггера \texttt{trg\_update\_season\_stats\_func}}
\label{fig:algo-trg-update}
\end{figure}
}

\newcommand{\flowchartTeamStats}{%
\begin{figure}[H]
\centering
\begin{tikzpicture}[
    fc base,
    node distance=1.2cm,
    startstop/.style={fc startstop},
    process/.style={fc process},
    arrow/.style={fc arrow}
]
\node (start) [startstop] {Начало построения представления \texttt{v\_team\_stats}};
\node (join)  [process, below of=start]
  {Соединение \texttt{game\_player\_stats} с \texttt{games}};
\node (role)  [process, below of=join]
  {Определение команды по \texttt{home\_team\_id} и \texttt{away\_team\_id}};
\node (group) [process, below of=role]
  {Группировка по \texttt{team\_id} и \texttt{season\_id}};
\node (agg)   [process, below of=group]
  {teamTotals := SUM(очки, подборы, \\ передачи, потери, броски)};
\node (avg)   [process, below of=agg]
  {Деление накопленных сумм на число матчей команды};
\node (stop)  [startstop, below of=avg] {Вывод средних командных показателей};

\draw [arrow] (start) -- (join);
\draw [arrow] (join)  -- (role);
\draw [arrow] (role)  -- (group);
\draw [arrow] (group) -- (agg);
\draw [arrow] (agg)   -- (avg);
\draw [arrow] (avg)   -- (stop);
\end{tikzpicture}
\caption{Схема алгоритма агрегирования командной статистики в представлении \texttt{v\_team\_stats}}
\label{fig:algo-team-stats}
\end{figure}
}

\newcommand{\flowchartUpdateSeasonStats}{%
\begin{figure}[H]
\centering
\begin{tikzpicture}[
    fc base,
    node distance=\flowchartUpdateVGap,
    startstop/.style={fc startstop},
    process/.style={fc process},
    predefined/.style={fc predefined},
    arrow/.style={fc arrow}
]
\node (start) [startstop] {Начало транзакции};
\node (pro1)  [process, below of=start]
  {Группировка по \texttt{player\_id}, \texttt{team\_id}: базовые агрегаты};
\node (pro2)  [predefined, below of=pro1]
  {Подзапрос TmFGA/TmFTA/TmTOV по (\texttt{game\_id}, \texttt{team\_id}) для USG};
\node (pro3)  [process, below of=pro2]
  {Двухпроходный расчёт лигового среднего (PER, BPM)};
\node (pro4)  [predefined, below of=pro3, yshift=\flowchartUpdateMidExtra]
  {Вызов \texttt{calculate\_efg}, \texttt{calculate\_ts}, \\ \texttt{calculate\_usg}, \texttt{calculate\_per}, \texttt{calculate\_bpm}};
\node (pro5)  [process, below of=pro4, yshift=\flowchartUpdateMidExtra]
  {Вставка с обновлением в таблицу \texttt{player\_season\_stats}};
\node (stop)  [startstop, below of=pro5] {Конец транзакции};

\draw [arrow] (start) -- (pro1);
\draw [arrow] (pro1)  -- (pro2);
\draw [arrow] (pro2)  -- (pro3);
\draw [arrow] (pro3)  -- (pro4);
\draw [arrow] (pro4)  -- (pro5);
\draw [arrow] (pro5)  -- (stop);
\end{tikzpicture}
\caption{Схема алгоритма хранимой процедуры \texttt{update\_season\_stats}}
\label{fig:algo-update}
\end{figure}
}
